\subsection{Controlled effect at fixed range and range interaction}
\label{sec:exp-identify}

We first ask whether interior allocation affects learned behavior when the sampled range is exactly fixed. We train paired $151.9$M models with $K=32$ rotary pairs at sequence length $256$. Within each of three seeds, architecture, initialization, token order, optimizer, schedule, and the $499{,}974{,}144$-token budget are identical. The geometric arm uses the FMRoPE training-base rule $b=L_{\mathrm{train}}$ \citep{oka2026fmrope}; the Cosh arm anchors its $\tau=4$ quantiles to the same endpoints. Only the $30$ interior exponents differ.

At the fixed training support, the mean Cosh-minus-geometric tail-NLL differences at lengths $256, 512, 1024$, and $2048$ are $+0.026, -0.281, -0.176, -0.146$. All three seeds improve at each extended length (Fig.~\ref{fig:spectral-budget-overview}). Interior allocation thus changes learned generalization without any scalar base shift. However, retargeting the log-frequency span to each evaluation length reverses the ordering to $+0.060, +0.227, +0.460$ across all three seeds. This reversal demonstrates that allocation cannot be judged independently of operating range: favorable fixed-support orderings do not survive arbitrary range dilation (Appendix~\ref{sec:identification-details}).

\subsection{Cross-shape factorial across configurations}
\label{sec:factorial}

A $50.9$M factorial varies two bases, two training lengths, three head dimensions, and three paired seeds, retaining endpoints throughout. At the reference Cosh strength, the weighted OOD objective improves in $7/12$ configurations; preassigned $1.25\times$ improves $10/12$, and matched exponential improves $9/12$. The aggregate reference-Cosh-minus-exponential contrast is small and uncertain ($+0.00074$ NLL, $95\%$ interval $[-0.006, 0.008]$; Table~\ref{tab:m4}). These results establish allocation effects across configurations while bounding claims about the unique superiority of any single shape.
